\textbf{Phần 1. Trắc nghiệm lựa chọn}
\setcounterpageref{ex}{0}

\begin{ex} %[2D10V51]
Số gồm $5$ trăm, $3$ chục và $8$ đơn vị được viết là
\choice
{$583$}
{$358$}
{$853$}
{\True $538$}
\loigiai{
Viết là $538$.
\textbf{Đáp án của câu hỏi là D}
}
\end{ex}

\begin{ex} %[2D10V51]
Số $706$ được đọc là
\choice
{Bảy mươi sáu}
{\True Bảy trăm linh sáu}
{Bảy trăm sáu mươi}
{Bảy không sáu}
\loigiai{
Số 706 đọc là Bảy trăm linh sáu.
\textbf{Đáp án của câu hỏi là B}
}
\end{ex}

\begin{ex} %[2D10V51]
Chữ số $4$ trong số $842$ có giá trị là
\choice
{$400$}
{$4$}
{\True $40$}
{$14$}
\loigiai{
Chữ số 4 ở hàng chục nên có giá trị là 40.
\textbf{Đáp án của câu hỏi là C}
}
\end{ex}

\begin{ex} %[2D10V51]
Số liền trước của số $300$ là
\choice
{\True $299$}
{$290$}
{$301$}
{$200$}
\loigiai{
$300 - 1 = 299$.
\textbf{Đáp án của câu hỏi là A}
}
\end{ex}

\begin{ex} %[2D10V51]
Từ ba chữ số $2, 5, 8$ (mỗi chữ số viết đúng một lần), có thể lập được bao nhiêu số có ba chữ số khác nhau?
\choice
{$4$ số}
{$5$ số}
{$3$ số}
{\True $6$ số}
\loigiai{
Có $3 \times 2 \times 1 = 6$ số: 258, 285, 528, 582, 825, 852.
\textbf{Đáp án của câu hỏi là D}
}
\end{ex}

\begin{ex} %[2D10V51]
Số có ba chữ số khác nhau lớn nhất được lập từ các chữ số $1, 9, 4$ là
\choice
{$914$}
{\True $941$}
{$491$}
{$194$}
\loigiai{
Xếp giảm dần: hàng trăm là 9, hàng chục là 4, hàng đơn vị là 1: số $941$.
\textbf{Đáp án của câu hỏi là B}
}
\end{ex}

\begin{ex} %[2D10V51]
Một số có chữ số hàng trăm là $6$, hàng đơn vị là $2$ và hàng chục gấp đôi hàng đơn vị. Số đó là
\choice
{$624$}
{$682$}
{\True $642$}
{$622$}
\loigiai{
Hàng đơn vị là 2, hàng chục gấp đôi là 4, hàng trăm là 6: số 642.
\textbf{Đáp án của câu hỏi là C}
}
\end{ex}

\begin{ex} %[2D10V51]
Số gồm $4$ trăm và $5$ đơn vị là số nào?
\choice
{\True $405$}
{$450$}
{$45$}
{$540$}
\loigiai{
Hàng chục bằng 0 nên viết là $405$.
\textbf{Đáp án của câu hỏi là A}
}
\end{ex}

\begin{ex} %[2D10V51]
Số lẻ liền sau số $539$ là
\choice
{$540$}
{$543$}
{$538$}
{\True $541$}
\loigiai{
Số lẻ liền sau là $539 + 2 = 541$.
\textbf{Đáp án của câu hỏi là D}
}
\end{ex}

\begin{ex} %[2D10V51]
Số trang của một cuốn sách có chữ số hàng trăm là $2$, hàng chục lớn hơn hàng trăm $3$ đơn vị, hàng đơn vị kém hàng chục $1$ đơn vị. Cuốn sách có bao nhiêu trang?
\choice
{$243$ trang}
{\True $254$ trang}
{$235$ trang}
{$255$ trang}
\loigiai{
Hàng trăm: 2; hàng chục: $2 + 3 = 5$; hàng đơn vị: $5 - 1 = 4$. Số trang là 254.
\textbf{Đáp án của câu hỏi là B}
}
\end{ex}

\begin{ex} %[2D10V51]
Có bao nhiêu số có ba chữ số có chữ số hàng trăm là $3$ và chữ số tận cùng là $0$?
\choice
{$9$ số}
{$8$ số}
{\True $10$ số}
{$11$ số}
\loigiai{
Các số dạng $3a0$ với $a \in \{0, 1, \dots, 9\}$. Có đúng 10 số.
\textbf{Đáp án của câu hỏi là C}
}
\end{ex}

\begin{ex} %[2D10V51]
Số bé nhất có ba chữ số mà tổng các chữ số bằng $3$ là
\choice
{\True $102$}
{$300$}
{$201$}
{$210$}
\loigiai{
Hàng trăm bé nhất là 1, hàng chục bé nhất là 0, hàng đơn vị là $3 - 1 - 0 = 2$: số 102.
\textbf{Đáp án của câu hỏi là A}
}
\end{ex}

\textbf{Phần 2. Trắc nghiệm đúng sai}
\setcounterpageref{ex}{0}

\begin{ex} %[2D10V51]
Cho số 825, xét tính đúng sai của các khẳng định sau:
\choiceTF
{\True Số 825 gồm 8 trăm, 2 chục và 5 đơn vị.}
{Chữ số 8 có giá trị là 80.}
{\True Số liền sau của 825 là 826.}
{\True Tổng các chữ số của 825 bằng 15.}
\loigiai{
\begin{itemize}
    \item Ý a) Đúng: Đúng cấu tạo thập phân.
    \item Ý b) Sai: Chữ số 8 ở hàng trăm nên có giá trị là 800.
    \item Ý c) Đúng: 825 + 1 = 826.
    \item Ý d) Đúng: 8 + 2 + 5 = 15.
\end{itemize}
\begin{itemize}
    \item a) \textbf{Đúng}.
    \item b) \textbf{Sai}.
    \item c) \textbf{Đúng}.
    \item d) \textbf{Đúng}.
\end{itemize}
}
\end{ex}

\begin{ex} %[2D10V51]
Xét các số có ba chữ số được lập từ các chữ số 0, 3, 7:
\choiceTF
{\True Số lớn nhất có ba chữ số khác nhau lập được là 730.}
{\True Số bé nhất có ba chữ số khác nhau lập được là 307.}
{Có thể lập được 6 số có ba chữ số khác nhau từ 3 chữ số trên.}
{\True Tổng của số lớn nhất và số bé nhất lập được là 1 037.}
\loigiai{
\begin{itemize}
    \item Ý a) Đúng: Đúng, xếp giảm dần: 730.
    \item Ý b) Đúng: Hàng trăm khác 0 nên chọn 3, tiếp theo là 0, 7. Đúng.
    \item Ý c) Sai: Vì chữ số 0 không đứng đầu nên chỉ lập được $2 \times 2 \times 1 = 4$ số.
    \item Ý d) Đúng: 730 + 307 = 1 037. Đúng.
\end{itemize}
\begin{itemize}
    \item a) \textbf{Đúng}.
    \item b) \textbf{Đúng}.
    \item c) \textbf{Sai}.
    \item d) \textbf{Đúng}.
\end{itemize}
}
\end{ex}

\textbf{Phần 3. Trả lời ngắn}
\setcounterpageref{ex}{0}

\begin{ex} %[2D10V51]
Viết số bé nhất có ba chữ số khác nhau được lập từ ba chữ số 0, 7, 3.
\shortans{307}
\loigiai{
\textbf{Đáp số:} $307$\\
Hàng trăm chọn 3, chục chọn 0, đơn vị chọn 7: số 307.
}
\end{ex}

\begin{ex} %[2D10V51]
Tìm một số có ba chữ số biết hàng trăm là 7, hàng chục là 0, xóa chữ số hàng trăm thì được số 9.
\shortans{709}
\loigiai{
\textbf{Đáp số:} $709$\\
Xóa hàng trăm còn lại 09 = 9. Vậy số đó là 709.
}
\end{ex}

\begin{ex} %[2D10V51]
Có bao nhiêu số có ba chữ số giống nhau?
\shortans{9}
\loigiai{
\textbf{Đáp số:} $9$\\
Đó là 111, 222, ..., 999. Có 9 số.
}
\end{ex}

\begin{ex} %[2D10V51]
Tìm số liền sau của số 799.
\shortans{800}
\loigiai{
\textbf{Đáp số:} $800$\\
799 + 1 = 800.
}
\end{ex}

\begin{ex} %[2D10V51]
Số có ba chữ số lớn nhất có các chữ số khác nhau là số nào?
\shortans{987}
\loigiai{
\textbf{Đáp số:} $987$\\
Đó là 987.
}
\end{ex}

\begin{ex} %[2D10V51]
Một số gồm 4 trăm, 0 chục và 9 đơn vị. Viết số đó.
\shortans{409}
\loigiai{
\textbf{Đáp số:} $409$\\
Viết là 409.
}
\end{ex}
